\documentclass{report}
\usepackage{amsfonts, amsmath}

\newcommand\N{{\mathbb N}}
\newcommand\Q{{\mathbb Q}}
\newcommand\R{{\mathbb R}}
\newcommand\Z{{\mathbb Z}}

\pagestyle{empty}
\begin{document}
\large
\centerline{\Huge \bf Analisi Matematica II}
\vskip.2cm
\centerline{\Huge  Secondo Appello (10-07-2002)}
\renewcommand{\thefootnote}{ } 
\footnote{\sl Avete 3 ore di tempo. Ogni esercizio
vale sei punti. La sufficienza si ottiene con un punteggio $\geq$
18. {\bf Solo le risposte chiaramente giustificate saranno prese in
considerazione.} {\scshape Elaborati illeggibili o confusi verranno ignorati.}
{\sffamily La
sintesi sar\`a  apprezzata.}}
\renewcommand{\thefootnote}{\arabic{footnote}} 

\vskip1cm
\begin{enumerate}
\item Sia $f:\R\to\R$ una funzione derivabile due volte e periodica di
periodo $2\pi$. Si 
consideri la funzione $g(x)=f(3x)$. Si dimostri che anche
$g$ \`e periodica di periodo $2\pi$. Si calcolino i coefficienti di
Fourier di $f$ in funzione di quelli di $g$. Si usi la relazione
trovata per dimostrare che se $f=g$ allora $f$ deve essere costante.
\item Sia $f:\R\to\R$ una funzione derivabile con derivata limitata da
$L$\  (cio\`e, $|f'(x)|\leq L$). Si calcoli il limite
\[
\lim_{n\to\infty}n\int_0^{\frac 1n}f(x)dx.
\]
\item Si calcoli il limite
\[
\lim_{x\to 0}\frac{2e^x-1-(x+1)^2}{x^3}.
\]
\item Sia $x:\R\to\R$ la coordinata di un punto che si muove su di una
retta secondo la legge
\[
\frac{dx}{dt}(t)=-\gamma x(t)+\gamma,\quad \gamma>0,
\]
e che si trova nel punto zero al tempo zero (cio\`e, $x(0)=0$). Si
studi il limite 
\[
\lim_{t\to\infty}x(t).
\]
\item  Si considerino le parabole $p_\lambda(x)=\lambda^{-4}x(\lambda-x)+1$
e sia $A(\lambda)$ l'area della regione $\{(x,y)\in\R^2\;|\;0\leq
x\leq \lambda;\; 0\leq y\leq p_\lambda(x)\}$. Si trovi il valore di
$\lambda>0$ per cui l'area $A(\lambda)$ \`e minima.
\item Si calcoli l'area di un settore circolare di angolo $\alpha\in(0,\pi/2)$ e
raggio $r$.
\end{enumerate}
\end{document}
